Obszary robocze, powszechnie zwane wirtualnymi pulpitami, pozwalają w łatwy sposób zarządzać i grupować programy dzięki czemu, przy dużej liczbie otwartych okien w łatwy sposób możemy zredukować zatłoczenie i usprawnić nawigację po naszym pulpicie. Dla przykładu, na jednym obszarze roboczym możemy przeglądać zasoby Internetu, a na drugim lub trzecim obszarze roboczym możemy mieć uruchomiony program do przeglądania multimediów oraz tworzyć dokument w programie biurowym.

Domyślnie Ubuntu daje dostęp do czterech obszarów roboczych, ale zaraz po instalacji ta opcja nie jest włączona. Aby móc korzystać z dodatkowych obszarów roboczych kliknij na ikonę Systemową \includegraphics{images/ikony_zasilanie.png} na panelu menu (w prawym górnym rogu) i wybierz \menu{{Ustawienia Systemu}>{Osobiste}>{Wygląd}>{Zachowanie}>{Uaktywnij obszary robocze}}.

\begin{wrapfigure}{l}{0.05\textwidth}
                \includegraphics[width=\linewidth]{images/ikony_obszary_robocze.png}
\end{wrapfigure}

Po zaznaczeniu dodatkowych obszarów roboczych, na pasku Launchera pojawi się dodatkowa ikona, która to służy do przełączania się pomiędzy obszarami roboczymi: \textcolor{ubuntu_orange}{Przełącznik obszarów roboczych}.
\clearpage
Domyślnie w Ubuntu masz do dyspozycji cztery wirtualne pulpity, ułożone w siatkę 2x2. Po naciśnięciu na ikonce obszarów roboczych na naszym ekranie pojawia się siatka czterech pulpitów. Pomiędzy obszarami roboczymi możesz się też przemieszczać za pomocą następujących skrótów klawiszowych:
\begin{itemize}
\item \keys{CTRL + ALT + \arrowkeyright} przemieszcza pulpit na obszar roboczy po prawej.
\item \keys{CTRL + ALT + \arrowkeyleft} przemieszcza pulpit na obszar roboczy po lewej.
\item \keys{CTRL + ALT + \arrowkeyup} przemieszcza pulpit na obszar roboczy na górze.
\item \keys{CTRL + ALT + \arrowkeydown} przemieszcza pulpit na obszar roboczy po na dole.
\end{itemize}